\textbf{COLUNI 2024 -- Questão 34}\par\medskip Observe a figura sobre uma cena comum nas grandes cidades.

\begin{center}
\includegraphics[width=0.48\linewidth]{static/simulados/coluni/figuras/2024/questao-34/figura-1.png}
\end{center}

\begin{center}Fonte: https://www.metropoles.com/conteudo-especial/mais-de-6-milhoes-de-pedidos-foram-entregues-pelo-projeto-ifood-pedal. Acesso em: 05 Ago. 2023.\end{center}

A atividade de entrega utilizando-se bicicletas tem se tornado mais comum em áreas urbanas. Em relação à geografia, esse trabalho pode estar diretamente ligado à dinâmica urbana e espacial das cidades. Marque a afirmativa que apresenta CORRETAMENTE as implicações geográficas que a profissão de entregadores ciclistas apresenta.
\begin{enumerate}
\item O trabalho dos \textit{bikers} pode ser influenciado pelo planejamento urbano, como a infraestrutura cicloviária:
      ciclovias, ciclofaixas, bicicletários, Centros de Atendimento ao Ciclista.
\item A modalidade de trabalho é realizada na grande maioria das vezes pelos aplicativos, que torna a situação
      ocupacional dos \textit{bikers} regulamentada e formal, com as garantias legais do trabalho.
\item A escolha da bicicleta como meio de transporte para entregas pode ter um impacto ambiental positivo,
      uma vez que não emite gases poluentes e contribui para a redução do tráfego de veículos motorizados.
\item Os entregadores ciclistas precisam otimizar suas rotas de entrega para serem eficientes e rápidos. Isso
      envolve conhecimento das ruas, das distâncias, do trânsito, da topografia, entre outros aspectos geográficos.
\end{enumerate}
